\section{Methods: fitting cycles, and the null they must beat}
\label{sec:methods}

We give the canonical cycles every reasonable opportunity to survive.
This section sketches the fitting arsenal and the null hypotheses they
are tested against. Formulae, parameter choices, implementation
references and unit-test pointers appear in Appendix~\ref{app:methods}.

\subsection{Cycle-fitting estimators}

A single cycle-detection method invites the rebuttal that the result is
an artefact of the method itself. We therefore commit, \emph{ex ante},
to four estimators drawn from four different families of statistical
assumption.

\paragraph{(a)~Christiano-Fitzgerald band-pass.}
The Christiano-Fitzgerald (2003) asymmetric filter
\citep{christiano_fitzgerald2003} is full-sample, optimal in mean
square for AR(1)-style processes, and dominates the symmetric
Baxter-King filter at the endpoints. The instantaneous phase is
extracted via the Hilbert transform.

\paragraph{(b)~Continuous Morlet wavelet.}
The continuous wavelet transform with a Morlet mother
(\(\omega_0=6.0\), \(\Delta j=0.125\)) is the standard tool for
time-frequency localisation \citep{torrence_compo1998,
grinsted_moore_jevrejeva2004} and is used here to compute the band-power
inside each canonical cycle band.

\paragraph{(c)~Markov-switching AR(1).}
Following \citet{hamilton1989}, a two-regime hidden-Markov model is
fitted by EM to the log-differenced series. A cycle is consistent with
periodic regime transitions whose dwell-time matches the canonical band.

\paragraph{(d)~PELT changepoints.}
The PELT algorithm of \citet{killick_fearnhead_eckley2012} detects
piecewise-constant mean shifts at \(O(n)\) cost under a BIC penalty,
providing an estimator orthogonal to the spectral and regime-switching
families.

\paragraph{(e)~Bry-Boschan turning points.}
The Bry-Boschan procedure \citep{bry_boschan1971} revived by
\citet{harding_pagan2002} identifies peaks and troughs from a censored
local-extremum rule (minimum phase duration: 5 quarters or 2 years),
giving a deterministic, non-parametric cycle estimator.

Estimators (a)--(d) are used as votes in the Gate~2 consensus
(Section~\ref{sec:criteria}); (e) is reported for cross-comparison.
The four-method consensus is intentionally over-determined: each
estimator embodies a generative hypothesis incompatible with the
others. If a cycle is real, all four should agree; if a cycle is an
artefact of one estimator, three should disagree.

\subsection{Two methods we do not use, and why}
\label{sec:methods_excluded}

\paragraph{Hodrick-Prescott filter.}
HP is excluded as a cycle-detection tool because
\citet{hamilton2018} shows that it introduces spurious cyclical
patterns even in pure random-walk data. Allowing HP into the consensus
would prejudice the test \emph{in favour} of the cycles---the opposite
of falsifiability.

\paragraph{Empirical Mode Decomposition (EMD/EEMD).}
EMD suffers from mode-mixing: a single intrinsic-mode function can
absorb multiple physical frequencies, making the decomposition
non-unique and non-falsifiable in our sense. The pre-registration
forbids non-unique decompositions.

\subsection{Null hypotheses: a dual-null surrogate test}
\label{sec:dual_null}

A statistically robust rejection of a cycle requires the band-power
to fail to dominate \emph{both} of the following nulls
\citep{theiler_et_al1992,schreiber_schmitz2000}.

\paragraph{Null~1: AR(1) bootstrap.}
We fit \(y_t = \mu + \phi(y_{t-1}-\mu) + \epsilon_t\) and simulate
\(B=1000\) independent realisations of identical length, mean,
variance and persistence. The empirical band-power is compared to
the 95th percentile of the surrogate band-powers.

\paragraph{Null~2: Phase-randomised (Theiler) surrogate.}
A Fast Fourier Transform of the series yields amplitudes
\(|\hat y_k|\) and phases \(\arg\hat y_k\); we draw a new phase
sequence uniformly on \([0,2\pi)\), enforce conjugate symmetry to
keep the inverse FFT real, and invert. The surrogate preserves the
empirical power spectrum but destroys cycle-level coherence.

\paragraph{The combined criterion.}
A cell is admitted at Gate~1 only if the empirical band-power
\emph{simultaneously} dominates Null~1 and Null~2 at \(\alpha=0.05\)
on \(B=1000\) surrogates. Failure to dominate \emph{either} null
yields a rejection. The dual-null is conservative by construction.

\subsection{Pre-registered bands}

The cycle bands are frozen in source code from the project's
inception\footnote{\path{ecowave/cycles/bands.py:CYCLE_BANDS} is
exported as an immutable \texttt{MappingProxyType}; a unit test
asserts immutability.}: Kitchin~\(3\text{--}5\)y
\citep{kitchin1923,wen2005,diebolt_doliger2008},
Juglar~\(7\text{--}11\)y \citep{juglar1862,schumpeter1939},
Kuznets~\(15\text{--}25\)y \citep{kuznets1930,korotayev_tsirel2010},
Kondratieff~\(40\text{--}60\)y
\citep{kondratieff1925,modelski1987,korotayev_tsirel2010}. We do
not adjust these bands \emph{post hoc}; the only band-related
robustness check we permit is the inclusion of a quarterly panel
to mitigate Nyquist concerns on Kitchin (Section~\ref{sec:data}).
